% =============================================================================
% Covariant entropy accounting for interacting ITSM sectors
% Version 0.7.1-conditional-draft
% This source preserves and supersedes the physical claims of v0.6.0.
% =============================================================================
\documentclass[aps,prd,reprint,superscriptaddress,nofootinbib,floatfix]{revtex4-2}

\usepackage[T1]{fontenc}
\usepackage[utf8]{inputenc}
\usepackage{amsmath,amssymb,amsthm}
\usepackage{booktabs}
\usepackage[hidelinks]{hyperref}

\newcommand{\Tthree}{T^{3}}
\newcommand{\dd}{\mathrm{d}}
\newcommand{\Qmp}{Q_{\mathrm{mp}}}
\newcommand{\Qsyn}{Q_{\mathrm{syn}}}

\hypersetup{
  pdftitle={Covariant entropy accounting for interacting matter-plenum sectors on compact spatial manifolds},
  pdfauthor={Brendon Boyd},
  pdfsubject={A bounded thermodynamic framework for interacting cosmological sectors},
  pdfkeywords={relativistic thermodynamics, interacting fluids, entropy production, cosmology, compact topology}
}

\begin{document}

\title{Covariant entropy accounting for interacting matter--plenum sectors
on compact spatial manifolds}

\author{Brendon Boyd}
\email{brendon.boyd@itsm-cosmology.org}
\affiliation{Independent Researcher, Burswood, Western Australia, Australia}
\thanks{%
\begin{tabular}{@{}l@{\ }l@{}}
ORCID: & \href{https://orcid.org/0009-0007-4177-2612}{0009-0007-4177-2612} \\
Repository: & \href{https://github.com/brendohxd/ITSM-Integrated-Toroidal-Syntropic-Model}{\texttt{github.com/brendohxd/ITSM-...}} \\
\end{tabular}%
}

\date{Version 0.7.1-conditional-draft --- \today}

\begin{abstract}
Energy exchange between cosmological sectors is not, by itself, entropy
production.  We derive the distinction for interacting perfect-fluid sectors
on a spacetime with compact spatial sections.  If
\(\nabla_\mu T_A^{\mu\nu}=Q_A^\nu\),
\(\nabla_\mu(n_Au^\mu)=n_A\Gamma_A\), and
\(q_A\equiv-u_\nu Q_A^\nu\), the Gibbs relation gives
\[
 \nabla_\mu S_A^\mu
 =\frac{q_A-\mu_A n_A\Gamma_A}{T_A}.
\]
Total stress conservation enforces \(\sum_A q_A=0\) in a common comoving
frame, but does not fix the sign of \(\sum_A\nabla_\mu S_A^\mu\).
For two particle-conserving sectors the result is
\(q(1/T_1-1/T_2)\), which is nonnegative only after a constitutive direction
for the transfer has been supplied.  Compactness removes net spatial boundary
flux; it does not create an external energy or entropy source.

We apply this accounting to the conditional R4C1 matter--plenum--reservoir
action.  Its action-derived currents and homogeneous balance equations support
finite reversible exchange, but the model supplies no temperatures, chemical
potentials, particle-production rates, dissipative transport matrix, or
coarse-graining prescription.  Consequently it does not yet derive positive
entropy production, heat-death avoidance, a thermodynamic stability Hessian,
or cosmic acceleration.  We formulate the constitutive, stability, and
observational conditions that a positive claim must satisfy.  The result is a
bounded thermodynamic framework and a set of no-go statements, not a completed
syntropic cosmology.
\end{abstract}

\keywords{relativistic thermodynamics; interacting fluids; entropy production;
cosmology; compact topology}

\maketitle

% =============================================================================
\section{Problem and scope}
\label{sec:scope}

Interacting-sector cosmologies commonly introduce currents through
\(\nabla_\mu T_A^{\mu\nu}=Q_A^\nu\).  Such equations are necessary for
consistent energy--momentum bookkeeping, but neither a nonzero \(Q_A^\nu\) nor
the identity \(\sum_AQ_A^\nu=0\) is a second-law statement.  Relativistic
non-equilibrium thermodynamics requires an entropy current and constitutive
relations in addition to stress-energy balance
\cite{israel_stewart,calvao_lima_waga,iosifidis2024}.

This distinction is especially important for models that use the label
\emph{syntropic}.  The name does not determine the direction of energy flow,
the sign of entropy production, or the stability of an equilibrium.  Those
are calculations.

Our objective is methodological: to separate what follows from covariant
balance identities from the additional physical input needed to assign
entropy, prove a second-law sign, establish equilibrium stability, or infer
cosmic acceleration.  Sections~\ref{sec:balances}--\ref{sec:topology} are
model independent within the stated local-equilibrium perfect-fluid
assumptions.  Section~\ref{sec:r4c1} then uses one conditional action as a
diagnostic application; the general entropy identity does not depend on that
model.

Table~\ref{tab:claims} records the resulting claim ladder.

\begin{table*}[t]
\caption{Status of thermodynamic claims in the present conditional model.}
\label{tab:claims}
\begin{ruledtabular}
\begin{tabular}{p{0.37\textwidth}p{0.16\textwidth}p{0.38\textwidth}}
Claim & Status & Evidence or missing input \\
\hline
Sector exchange currents and total balance & Conditional derivation &
R4C1 action variation and homogeneous on-shell control \\
Finite, reversible exchange on a sampled background & Locally checked &
Both interface currents are nonzero and the reservoir transfer reverses sign \\
Positive total entropy production & Not derived &
No \(T_A,\mu_A,\Gamma_A\), dissipative kernel, or coarse graining \\
Topology supplies external energy/entropy inflow & Excluded &
\(\int_{\Tthree}\sqrt{\gamma}\,D_iJ^i=0\) for smooth periodic \(J^i\) \\
Thermodynamic stability & Open &
No constrained entropy Hessian or susceptibility matrix has been computed \\
Acceleration from exchange & Not derived &
Requires the total pressure and a solved background; \(q_A\) cancels in the
summed continuity equation \\
Avoidance of heat death & Not established &
Requires a global state, asymptotics, and entropy definition beyond the
present classical reversible action \\
\end{tabular}
\end{ruledtabular}
\end{table*}

% =============================================================================
\section{Covariant sector balances}
\label{sec:balances}

Consider sectors \(A=1,\ldots,N\) with symmetric stresses
\(T_A^{\mu\nu}\) and transfer currents
\begin{equation}
 \nabla_\mu T_A^{\mu\nu}=Q_A^\nu,
 \qquad
 \sum_A Q_A^\nu=0.
 \label{eq:sector-balance}
\end{equation}
The second relation follows from total diffeomorphism invariance after all
field equations are imposed; it should not be inserted as a substitute for
varying an incompletely specified action.

Let \(u^\mu u_\mu=-1\), \(h^\mu{}_\nu=\delta^\mu{}_\nu+u^\mu u_\nu\),
and decompose each current as
\begin{equation}
 q_A\equiv-u_\nu Q_A^\nu,
 \qquad
 f_A^\mu\equiv h^\mu{}_\nu Q_A^\nu,
 \qquad
 Q_A^\mu=q_Au^\mu+f_A^\mu .
 \label{eq:current-decomp}
\end{equation}
For a perfect fluid,
\begin{equation}
 T_A^{\mu\nu}=(\rho_A+p_A)u^\mu u^\nu+p_Ag^{\mu\nu},
\end{equation}
the timelike projection of Eq.~\eqref{eq:sector-balance} is
\begin{equation}
 \dot\rho_A+\theta(\rho_A+p_A)=q_A,
 \qquad
 \theta\equiv\nabla_\mu u^\mu.
 \label{eq:energy-balance}
\end{equation}
In an FLRW chart, \(\theta=3H\).  Equation~\eqref{eq:energy-balance} fixes an
energy transfer convention: \(q_A>0\) means energy enters sector \(A\) in the
chosen comoving frame.  The spatial projection gives its force density and
must be retained in perturbative or non-comoving applications.

Summing Eq.~\eqref{eq:energy-balance} gives
\begin{equation}
 \dot\rho_{\rm tot}+3H(\rho_{\rm tot}+p_{\rm tot})=0,
 \label{eq:total-continuity}
\end{equation}
for a common FLRW velocity.  This identity is conservation, not an equation
for entropy.

% =============================================================================
\section{Gibbs relation and entropy production}
\label{sec:gibbs}

Introduce a particle current and specific entropy for each sector,
\begin{equation}
 N_A^\mu=n_Au^\mu,
 \qquad
 \nabla_\mu N_A^\mu=n_A\Gamma_A,
 \qquad
 S_A^\mu=n_A\sigma_Au^\mu .
 \label{eq:number-entropy}
\end{equation}
Here \(\Gamma_A\) is a fractional particle-production rate.  It is independent
information: an energy-transfer current does not determine it.

For local equilibrium the Gibbs relation per particle is
\begin{equation}
 T_A\dd\sigma_A
 =\dd\!\left(\frac{\rho_A}{n_A}\right)
 +p_A\,\dd\!\left(\frac{1}{n_A}\right).
 \label{eq:gibbs}
\end{equation}
Taking the derivative along \(u^\mu\) and using
Eqs.~\eqref{eq:energy-balance} and \eqref{eq:number-entropy} yields
\begin{equation}
 n_AT_A\dot\sigma_A
 =q_A-(\rho_A+p_A)\Gamma_A .
 \label{eq:specific-entropy}
\end{equation}
The Euler relation
\(\rho_A+p_A=T_An_A\sigma_A+\mu_An_A\) then gives the entropy-current
divergence
\begin{equation}
 \boxed{\;
 \nabla_\mu S_A^\mu
 =\frac{q_A-\mu_An_A\Gamma_A}{T_A}
 \;}
 \label{eq:entropy-production}
\end{equation}
for the idealized perfect-fluid sector.  Heat flux, viscosity, diffusion,
chemical reactions, and relaxation add further flux--force products; causal
closures require more than the algebraic first-order form
\cite{israel_stewart}.

\subsection{Two particle-conserving sectors}

Set \(\Gamma_1=\Gamma_2=0\), \(q_1=q\), and \(q_2=-q\).  Then
\begin{equation}
 \nabla_\mu(S_1^\mu+S_2^\mu)
 =q\left(\frac{1}{T_1}-\frac{1}{T_2}\right).
 \label{eq:two-sector-entropy}
\end{equation}
The right-hand side can be positive, zero, or negative for an arbitrarily
assigned \(q\).  A simple Onsager closure
\begin{equation}
 q=L\left(\frac{1}{T_1}-\frac{1}{T_2}\right),
 \qquad L\geq0,
 \label{eq:onsager}
\end{equation}
gives
\begin{equation}
 \nabla_\mu S_{\rm tot}^\mu
 =L\left(\frac{1}{T_1}-\frac{1}{T_2}\right)^2\geq0.
 \label{eq:onsager-positive}
\end{equation}
Thus nonnegative entropy production is a property of a constitutive law and
its coefficients, not of \(\sum_AQ_A^\nu=0\).

\subsection{Particle creation}

When \(\Gamma_A\neq0\), chemical potentials and particle source terms enter
Eq.~\eqref{eq:entropy-production}.  Constant specific entropy
\(\dot\sigma_A=0\) imposes
\begin{equation}
 q_A=(\rho_A+p_A)\Gamma_A,
 \label{eq:adiabatic-creation}
\end{equation}
but this is a constitutive restriction rather than a consequence of the
continuity equation.  Cosmological matter-creation thermodynamics has long
made this distinction explicit \cite{prigogine1988,calvao_lima_waga}.

% =============================================================================
\section{Compact topology: what vanishes and what does not}
\label{sec:topology}

Let each spatial slice be a smooth flat \(\Tthree\) with periodic fields and
induced metric \(\gamma_{ij}\).  For any smooth periodic spatial vector,
\begin{equation}
 \int_{\Tthree}\dd^3x\,\sqrt{\gamma}\,D_iJ^i=0.
 \label{eq:compact-divergence}
\end{equation}
This is the divergence theorem on a compact manifold without boundary.  It
forbids interpreting a pure divergence as net inflow from an exterior.

For a homogeneous entropy current, the integrated entropy balance is
\begin{equation}
 \frac{\dd}{\dd t}
 \left(a^3V_0\sum_A n_A\sigma_A\right)
 =a^3V_0\sum_A
 \frac{q_A-\mu_An_A\Gamma_A}{T_A},
 \label{eq:integrated-entropy}
\end{equation}
where \(V_0\) is the coordinate volume of the fundamental domain.  Compactness
removes a boundary term from this equation; it does not set the right-hand
side to zero and does not choose its sign.  Internal energy transfer and local
entropy production remain possible.

The same boundary fact also makes zero modes and global constraints important.
A spatially constant mode is not fixed by Eq.~\eqref{eq:compact-divergence};
its state, charge, and thermodynamic interpretation require separate input.

\subsection{Provenance of the topology correction}

Earlier ITSM working notes asserted that compact \(\Tthree\) topology requires
an inward energy stream and that matter--plenum exchange exactly cancels
entropy growth.  Both statements are withdrawn.  Equation~\eqref{eq:compact-divergence}
excludes the proposed exterior boundary flux, while
Eq.~\eqref{eq:entropy-production} shows why total energy conservation alone
cannot cancel entropy production.  The previous v0.6.0 source and PDF remain
preserved as provenance; the present manuscript supersedes their
thermodynamic claims.

% =============================================================================
\section{Application to the conditional R4C1 action}
\label{sec:r4c1}

\subsection{Why R4C1 is used as a stress test}

The general accounting boundary becomes operational when applied to a model
whose interface currents are independently obtained from an action.  R4C1 is
useful here because it contains several dynamical sectors and nonzero,
action-derived exchange currents, including a reservoir transfer that changes
sign on the registered trajectory, while supplying no dissipative
constitutive closure.  It is therefore sufficiently specified to test
conservation and finite reversible exchange, but deliberately insufficient to
infer irreversible entropy production.  This application diagnoses the gap
between those statements; it is not evidence that R4C1 is a completed or
observationally viable cosmology.

\subsection{Action-derived currents}

The R4C1 candidate is a closed, reversible, classical field action on
\(I\times\Tthree\).  It contains a conformally coupled dust sector, a complex
amplitude/phase sector with an independent force chart, a constrained timelike
frame, and a real reservoir field.  The baseline stress split gives the
action-derived interface currents
\begin{align}
 \Qmp^\nu&=\beta T_m\nabla^\nu\psi, \\
 \Qsyn^\nu&=-\frac{g_r}{2}s\,r\,\nabla^\nu r,
 \qquad s=u^2+v^2 .
 \label{eq:r4c1-currents}
\end{align}
On all field equations the corresponding sector identities are
\begin{align}
 \nabla_\mu T_m^{\mu\nu}&=\Qmp^\nu,\\
 \nabla_\mu T_P^{\mu\nu}&=-\Qmp^\nu+\Qsyn^\nu,\\
 \nabla_\mu T_R^{\mu\nu}&=-\Qsyn^\nu.
 \label{eq:r4c1-balances}
\end{align}
Their sum vanishes.

On the registered homogeneous background these reduce to
\begin{align}
 \dot\rho_m+3H\rho_m&=\Qmp^0,\\
 \dot\rho_P+3H(\rho_P+p_P)&=-\Qmp^0+\Qsyn^0,\\
 \dot\rho_R+3H(\rho_R+p_R)&=-\Qsyn^0.
 \label{eq:r4c1-frw}
\end{align}
The numerical control has nonzero interface currents and a finest independently
differenced sector-balance residual of \(7.96\times10^{-8}\).  The reservoir
transfer changes sign over the sampled trajectory.  This is evidence for
finite reversible energy exchange in the conditional model, not for an
irreversible pump.

\subsection{Why no entropy-production claim follows}

The R4C1 action does not currently define:
\begin{enumerate}
\item sector temperatures \(T_m,T_P,T_R\);
\item chemical potentials or particle currents for all sectors;
\item nonzero production rates \(\Gamma_A\);
\item a dissipative response matrix such as \(L_{AB}\);
\item a coarse-graining map from the reversible field dynamics to macroscopic
entropy; or
\item a quantum state and renormalized entropy for the compact modes.
\end{enumerate}
Consequently Eq.~\eqref{eq:entropy-production} cannot yet be evaluated.

There is a second caution.  Assigning a fraction \(\eta\) of the portal
potential to a different sector changes the named reservoir current by
\(\eta\nabla^\nu W\) while leaving the total action and stress unchanged.
A sector entropy source built from a split-dependent current is not an
observable until the subsystem convention and coarse graining are physically
fixed.

The action also conserves its full global \(U(1)\) current.  A nonzero
condensate-number source has not been derived.  Energy exchange and charge
creation must not be conflated.

% =============================================================================
\section{Thermodynamic stability is a constrained Hessian problem}
\label{sec:stability}

The second law and stability are different tests.  A nonnegative instantaneous
entropy-production rate does not prove that an equilibrium is a local maximum
of entropy.

Let \(X^a\) denote independent extensive variables after all conserved totals
and constraints have been imposed.  Local thermodynamic stability requires
\begin{equation}
 H^{(S)}_{ab}
 \equiv
 \frac{\partial^2S_{\rm tot}}{\partial X^a\partial X^b}
 \preceq0
 \label{eq:entropy-hessian}
\end{equation}
on the allowed tangent space.  Equivalently, an energy representation requires
the appropriate susceptibility matrix to be positive semidefinite.

For two subsystems exchanging only energy at fixed total energy,
\begin{equation}
 \frac{\dd^2 S_{\rm tot}}{\dd E_1^2}
 =-\frac{1}{T_1^2C_1}
  -\frac{1}{T_2^2C_2}.
 \label{eq:heat-capacity}
\end{equation}
Positive heat capacities are sufficient for stability along this one transfer
direction, but a coupled matter--plenum system can have additional density,
charge, amplitude, phase, and modulus directions.  Each physical eigenvalue of
the constrained Hessian must be reported over a registered domain.

The current R4C1 quadratic-action work concerns dynamical kinetic, gradient,
constraint, and mass matrices.  Those are necessary for field stability but
are not interchangeable with Eq.~\eqref{eq:entropy-hessian}.  No numerical
thermodynamic Hessian eigenvalues are claimed here.

% =============================================================================
\section{Exchange does not by itself cause acceleration}
\label{sec:acceleration}

In Einstein gravity an FLRW background obeys
\begin{equation}
 \frac{\ddot a}{a}
 =-\frac{1}{6M_P^2}
 \sum_A(\rho_A+3p_A).
 \label{eq:acceleration}
\end{equation}
The transfer terms cancel from the summed continuity equation
\eqref{eq:total-continuity}.  Exchange may alter the sector partition and hence
the pressures, but acceleration requires a solved stress history satisfying
Eq.~\eqref{eq:acceleration}.  Entropy production, an acceleration transition,
and thermodynamic stability are therefore three distinct questions.

This separation has become observationally concrete.  A recent background
analysis of thermogravity found that acceleration generated solely by universal
energy non-conservation is strongly disfavored by supernovae plus DESI BAO,
with a reported \(\Delta\chi^2\simeq+51\) relative to flat
\(\Lambda\)CDM \cite{lee2026}.  R4C1 is not that model: its total stress is
conserved.  The comparison nevertheless shows why an exchange mechanism
cannot be promoted from formal consistency to cosmological success without
deriving \(H(z)\), the ruler calibration, and the observational dictionary.

% =============================================================================
\section{Methodological conditions for syntropic cosmological claims}
\label{sec:programme}

To avoid conflating conservation identities with thermodynamic closure, a
model in the interacting-sector class considered here that advances a positive
syntropic claim must provide the following model-specific ingredients and
tests:
\begin{enumerate}
\item a declared physical subsystem split and equations of state
\(\rho_A(n_A,\sigma_A,\ldots)\);
\item temperatures, chemical potentials, entropy currents, and particle
currents derived from that state space;
\item action-derived or microphysically derived \(\Gamma_A\) where particle
creation is claimed;
\item a causal constitutive closure for heat, diffusion, viscosity, chemical
reactions, and reservoir transfer, with a positive entropy-production form;
\item the constrained entropy Hessian and all eigenvalues over a registered
domain;
\item compatibility between thermodynamic and quadratic-action stability,
including constraint singularities.
\end{enumerate}

These conditions are necessary within the local-equilibrium framework used
here; they are not jointly sufficient, and alternative thermodynamic
formalisms must supply equivalent physical closure.  If a claim extends to
cosmic acceleration or observational viability, it additionally requires the
full action-derived \(H(z)\), \(q(z)\), sound horizon, BBN history, and
perturbations, followed by common-prior likelihood comparisons.  Independent
external review remains a publication requirement, not a substitute for any
physical item above.

A failure of any positivity, causality, uniqueness, or observational test must
remain visible.  Conversely, passing a conservation residual or a Hessian
check alone would not establish the rest of the chain.

% =============================================================================
\section{Conclusions}
\label{sec:conclusions}

The covariant entropy identity
\eqref{eq:entropy-production} supplies the missing accounting boundary for
interacting ITSM sectors.  It establishes four points.

First, total stress conservation and entropy production are logically
independent.  Second, compact \(\Tthree\) topology removes external boundary
flux rather than supplying it.  Third, R4C1 presently realizes finite,
action-derived, reversible exchange but does not contain the constitutive data
needed to prove positive entropy production or particle creation.  Fourth,
thermodynamic stability, dynamical field stability, and cosmic acceleration
are separate calculations.

The resulting audit sequence is portable across interacting-sector models:
derive the currents, specify the thermodynamic state space and constitutive
closure, evaluate entropy production, test the constrained thermodynamic and
field-stability spectra, and only then confront cosmological observables.
R4C1 currently reaches the first step and a bounded reversible-exchange
control.

Within that scope, this work provides a covariant framework in which proposals
for entropy reduction, heat-death avoidance, or exchange-driven acceleration
can be made falsifiable without confusing conservation identities with
thermodynamic closure.  It establishes none of those positive mechanisms in
the present conditional model.

\section*{Reproducibility and provenance}

The conditional action is frozen in
\texttt{RES001\_R4C1\_ACTION\_FREEZE\_2026-09-25.md}
(SHA-256 prefix \texttt{81bfc2e4f17b}).  The first-variation report containing
Eq.~\eqref{eq:r4c1-currents} has prefix \texttt{4834ea7517fd}; the interacting
background report containing Eq.~\eqref{eq:r4c1-frw} has prefix
\texttt{48b49fddd906}.  Full digests, contracts, executables, and JSON receipts
are retained in the companion repository.  Hashes establish artifact identity,
not physical truth or independent review.

\section*{Declaration on the use of artificial-intelligence tools}

Artificial-intelligence-assisted tools were used for editorial drafting,
source-code assistance, and checks of symbolic and numerical consistency.
The author independently reviewed the manuscript, equations, cited sources,
and generated artifacts and accepts full responsibility for the content.

\begin{acknowledgments}
The author thanks the reviewers and collaborators who insisted on separating
energy conservation, entropy production, stability, and acceleration.
Remaining errors are the author's.
\end{acknowledgments}

\bibliography{thermodynamics_v0.7.1}

\end{document}
